\documentclass[11pt]{article}
\usepackage{fontspec,amsmath,amssymb,amsthm,hyperref}
\usepackage[margin=1in]{geometry}
\setlength{\emergencystretch}{3em}
\title{An explicit obstruction to the isotropic GLM Hessian estimate}
\author{LangevinLean mathematical exposition}
\date{}
\begin{document}\maketitle
\noindent Source module: \texttt{AverageLMC/GLMObstruction.lean}. Compile this proof document with XeLaTeX or LuaLaTeX. The accompanying HTML pairs the mathematical argument with the exact Lean statements and records the scope of the formalization.


The isotropic row of the paper's GLM proposition asserts a Hessian operator-norm Lipschitz bound of order n, with a multiplicative constant depending only on the component derivative bound C\_g. This module checks a permitted family with one component g\_1(t)=t squared minus sin(t) and every other component equal to t squared. It identifies the actual finite-sum potential, its gradient and its Hessian, and proves that every Hessian Lipschitz constant must be at least the cube of the first design vector's norm. The Gaussian probability argument is proved separately. These are checked obstruction ingredients, not proofs of the paper's convergence theorems.



The fixed component family is \(g_1(t)=t^2-\sin t\), \(g_i(t)=t^2\) for \(i\ne1\). Its actual Hessian is \[H(x)=Q+\sin\langle a,x\rangle\,aa^\top,\qquad Q=2\sum_iX_iX_i^\top.\] Every operator norm Lipschitz bound M for H satisfies \(\|a\|^3\le M\).


\section{The oscillating component and its actual derivatives}


Define g(t)=t squared minus sin(t), g'(t)=2t minus cos(t), and g''(t)=2 plus sin(t). Differentiate these formulas using the power and trigonometric derivative rules; the third derivative is cos(t). These are derivatives of the same component, rather than independently supplied derivative fields. Polynomial and sine functions are smooth on the whole real line, so this choice satisfies the component regularity in the manuscript.



\noindent Lean declarations: \texttt{oscillatingComponent}, \texttt{oscillatingComponentGradient}, \texttt{oscillatingComponentHessian}, \texttt{hasDerivAt\_oscillatingComponent}, \texttt{hasDerivAt\_oscillatingComponentGradient}, \texttt{hasDerivAt\_oscillatingComponentHessian}, \texttt{oscillatingComponent\_contDiff}.

\section{A fixed derivative bound independent of dimension}


The inequality sin(t) at least minus one gives g''(t) at least one. Therefore g'' is nonnegative and its absolute value is itself. The bounds sin(t) at most one and absolute cos(t) at most one yield absolute g''(t) plus absolute g'''(t) at most four, uniformly over t. The component selector uses this oscillating function at the distinguished index and the quadratic at every other index. Each quadratic has second derivative two and third derivative zero. The module proves smoothness, convexity and the bound on actual second and third derivatives for every selected component. Thus the same fixed C\_g=4 applies to every component in every dimension.



\noindent Lean declarations: \texttt{oscillatingComponentHessian\_lower}, \texttt{oscillatingComponent\_derivative\_bound}, \texttt{oscillatingComponent\_convex}, \texttt{counterexampleComponent}, \texttt{counterexampleComponent\_contDiff}, \texttt{counterexampleComponent\_convex}, \texttt{counterexampleComponent\_derivative\_bound}.

\section{Write the finite-sum potential and its score explicitly}


For a finite design family X\_i and a distinguished index i\_0, the potential is the sum of the squared inner products with X\_i, minus one copy of sin of the inner product with X\_i0. Expanding the selector and summing the unique sine correction proves that this is exactly the GLM sum with the certified component family. Its smoothness on the whole state space follows from finite sums and composition of smooth polynomials, linear functionals and sine. Its constant quadratic Hessian Q is twice the sum of the rank-one operators X\_i X\_i transpose. The score is Qx minus cos of the inner product with a times a, where a=X\_i0. Differentiate each squared linear functional, sum the resulting continuous linear derivatives, and subtract the derivative of the single sine term. The result is the inner-product functional represented by this explicit score.



\noindent Lean declarations: \texttt{counterexamplePotential}, \texttt{counterexamplePotential\_eq\_sum}, \texttt{counterexamplePotential\_contDiff}, \texttt{quadraticHessian}, \texttt{glmGradient}, \texttt{hasFDerivAt\_counterexamplePotential}.

\section{Differentiate the score to obtain the Hessian}


Differentiate Qx using the constant linear map Q. The derivative of the negative cosine term is sin of the inner product with a times the rank-one operator aa transpose. Thus the Hessian is H(x)=Q plus sin of the inner product with a times aa transpose. This formula is proved as the Frechet derivative of the explicit score. In a complete real inner-product space, the previously proved derivative of the potential identifies this score as its actual gradient. Taking that gradient's derivative then identifies H as the actual Hessian. The operator norm used for these continuous linear maps is exactly the Euclidean operator norm relevant to the manuscript's M\_2.



\noindent Lean declarations: \texttt{glmHessian}, \texttt{hasFDerivAt\_glmGradient}, \texttt{gradient\_counterexamplePotential}, \texttt{fderiv\_gradient\_counterexamplePotential}.

\section{Differentiate the Hessian along the design direction}


Restrict the Hessian to the path x=t a. Its sine coefficient becomes sin(t times the squared norm of a). At t=0 this scalar function has derivative equal to the squared norm of a. Consequently the derivative of H(t a) at zero is the squared norm of a times the rank-one operator aa transpose. The rank-one operator has norm equal to the squared norm of a, so the Hessian path derivative has norm equal to the fourth power of the norm of a.



\noindent Lean declarations: \texttt{hasDerivAt\_glmHessian\_line}.

\section{Every Hessian Lipschitz constant has cubic growth}


Suppose H is M-Lipschitz in operator norm, with M nonnegative. Along x=t a, the path is M times the norm of a Lipschitz, because its state displacement has norm absolute t times the norm of a. A derivative of a Lipschitz path has norm at most its Lipschitz constant. The preceding derivative calculation therefore gives the fourth power of the norm of a at most M times the norm of a. If a is nonzero, divide by its positive norm to obtain the cube of its norm at most M. If a is zero, the result follows from M nonnegative. With n=p and a standard Gaussian first design vector, the separate Gaussian module shows why this necessary cubic lower bound prevents a uniform order-n bound with probability one half.



\noindent Lean declarations: \texttt{glmHessian\_lipschitz\_lower}, \texttt{counterexamplePotential\_hessian\_lipschitz\_lower}.

\end{document}
